\documentclass[11pt,a4paper]{article}
\usepackage[utf8]{inputenc}
\usepackage{cmap}
\usepackage[T1]{fontenc}
\usepackage{lmodern}
\usepackage[spanish,es-noquoting]{babel}
\usepackage[margin=2.3cm]{geometry}
\usepackage{booktabs}
\usepackage{parskip}
\usepackage{amsmath}
\usepackage{caption}
\renewcommand{\arraystretch}{1.15}
\captionsetup{font=small,labelfont=bf}
\setlength{\parskip}{0.55em}

\begin{document}

\begin{center}
{\large\bfseries Informe de cierre --- frente de notas de rescate}\\[6pt]
{\small 28 de septiembre de 2026 \quad\textbullet\quad Para: Prof.\ Cristian Cappo \quad\textbullet\quad
De: Romina Alfonzo y Carlos Urdapilleta}\\[4pt]
{\footnotesize\itshape Detección de familias de ransomware en base a archivos encriptados
y notas de rescate}
\end{center}

\vspace{0.3em}
\hrule
\vspace{0.9em}

El frente de notas queda cerrado. El sistema \textbf{acierta la familia en el 81,2\,\% de las notas}
y alcanza macro-F1 \textbf{0,7417} medido en el escenario \emph{más exigente} de todos: una plantilla
de nota que nunca vio. Para dimensionarlo: el esquema de catálogo cerrado que emplea la literatura
previa, medido sobre este mismo corpus y con el mismo método, da 0,7767. \textbf{Nuestro resultado en
el caso difícil está prácticamente al nivel del que se reporta en el caso fácil.} Y en la situación
más común en la práctica --- que la nota se parezca a alguna ya conocida de su familia --- el sistema
\textbf{acierta el 98,9\,\%}.

Durante el cierre detectamos y corregimos además un defecto en nuestra propia partición de
evaluación, que es lo que explica la diferencia con las cifras que le habíamos informado antes:
\textbf{el sistema no cambió; cambió la forma de medirlo}. Este informe dice qué se construyó, qué da
en cada escenario, dónde falla y qué limitaciones declaramos. El detalle completo, con bases y semillas, está en
\texttt{ESTADO\_TESIS.md} y en el Capítulo~4. El frente de archivos cifrados se informa por
separado, como corresponde a dos sistemas independientes.

\section*{1. Qué se construyó}

Un identificador de familia que, ante una nota de rescate, decide en tres pasos, igual que un
analista:

\begin{enumerate}\setlength{\itemsep}{0.15em}
  \item \textbf{Indicadores propios de la nota}: correos de contacto, direcciones de pago, URLs y
        formato del identificador de víctima. El diccionario valor $\rightarrow$ familia se arma
        \emph{solo} con el conjunto de entrenamiento, y se descarta todo valor que en entrenamiento
        aparezca en más de una familia.
  \item \textbf{Nombre genuino del archivo} de la nota, cuando está auditado.
  \item \textbf{El texto}, con TF-IDF sobre caracteres y palabras y un clasificador lineal de
        margen máximo.
\end{enumerate}

Los dos primeros pasos responden solo por unanimidad; ante conflicto o ausencia de coincidencia,
decide el texto.

\textbf{Corpus:} 149 notas auténticas de 30 familias --- las mismas 30 de NapierOne, que emparejan
los dos frentes del trabajo --- agrupadas en 99 plantillas por casi-duplicado. \textbf{Cero notas
sin fuente citable}, y el manifiesto cuadra uno a uno con el disco.

\section*{2. Cómo se evalúa, y una corrección que hicimos en el camino}

La pregunta que importa en un despliegue real no es si el sistema reconoce una nota que ya vio,
sino si reconoce una \textbf{plantilla nunca vista}. Por eso el corpus se parte \emph{por
plantilla}: todas las notas de una plantilla caen del mismo lado, de modo que el texto de prueba
nunca estuvo en entrenamiento.

Durante el cierre detectamos que nuestro repartidor de plantillas, si bien respetaba esa garantía,
\textbf{no aseguraba que cada familia tuviera al menos un ejemplo en entrenamiento en cada
pliegue}. Dejaba 3,86 familias por pliegue sin material alguno, que sacaban F1 = 0 sin que el
método hubiera fallado, y esos ceros entraban al promedio macro.

El reparto corregido distribuye las plantillas \emph{dentro} de cada familia. Mantiene todo lo
demás igual, y se verificó:

\begin{center}
\small
\begin{tabular}{lcc}
\toprule
\textbf{Control} & \textbf{Reparto anterior} & \textbf{Reparto corregido} \\
\midrule
Plantilla en entrenamiento y prueba a la vez & 0 & 0 \\
Plantillas de entrenamiento por pliegue & 49,5 & 49,5 \\
Familias sin entrenamiento por pliegue & 3,86 & 1,00 \\
\bottomrule
\end{tabular}
\end{center}

De esas 3,86 familias, \textbf{1,00 es inevitable} --- las de plantilla única, ausentes en uno de
los dos pliegues --- y \textbf{2,86 eran defecto del reparto}. El tamaño de entrenamiento es
idéntico: la mejora no viene de entrenar con más datos.

El preregistro de esta corrección, con siete predicciones falsificables y una puerta de entrada que
aborta si no reproduce el evaluador anterior, quedó \textbf{commiteado en el repositorio antes de
correr el experimento}, de modo que la marca temporal es verificable.

\section*{3. Resultados}

\subsection*{Qué tan bien funciona, según el escenario}

Una misma herramienta da números muy distintos según qué se le pida. Medimos los tres escenarios
\emph{sobre el mismo corpus y con el mismo método}, de modo que lo único que cambia entre filas es la
pregunta que se responde:

\begin{center}
\small
\begin{tabular}{llc}
\toprule
\textbf{Escenario} & \textbf{Qué mide} & \textbf{macro-F1} \\
\midrule
Catálogo cerrado & buscar la nota más parecida del catálogo; & 0,7767 \\
(esquema de la literatura previa) & es lo que reporta el trabajo previo & \\[3pt]
Nota nueva de una plantilla ya conocida & reconocer una variante de algo visto & 0,7889 \\[3pt]
\textbf{Plantilla nunca vista} & \textbf{el caso de despliegue real} & \textbf{0,7417} \\
\textbf{(sistema completo)} & & \\
\bottomrule
\end{tabular}
\end{center}

La lectura importa: \textbf{nuestro número en el escenario más duro (0,7417) es comparable al que el
esquema previo obtiene en el más benévolo (0,7767)}, y lo único que cambia entre filas es \emph{el
protocolo de evaluación}, no el método. Esa comparación es, en sí misma, uno de los aportes del
trabajo.

Dentro del escenario exigente, el desempeño depende de si la nota se parece a algo ya conocido:

\begin{center}
\small
\begin{tabular}{lcc}
\toprule
\textbf{Situación de la nota} & \textbf{Cuántas} & \textbf{Acierto} \\
\midrule
Se parece a una nota conocida de su familia & 54 (36,2\,\%) & \textbf{0,9891} \\
No se parece a ninguna conocida & 91 (63,8\,\%) & 0,7434 \\
\midrule
\textbf{Global} & 145 & \textbf{0,8123} \\
\bottomrule
\end{tabular}
\end{center}

Es decir que el sistema \textbf{resuelve casi sin error el caso frecuente} y sigue siendo útil en el
caso en que no tiene ningún antecedente parecido, que es justamente donde una herramienta de catálogo
falla por completo.

\subsection*{La cifra de cabecera}

Corpus de 149 notas, 30 familias, 50 semillas, plantilla nunca vista.

\begin{center}
\small
\begin{tabular}{lcccccc}
\toprule
 & \textbf{macro-F1} & \textbf{IC 95\,\%} & \textbf{28 eval.} & \textbf{Exact.} & \textbf{Bal.} & \textbf{MCC} \\
\midrule
\textbf{Sistema completo (cascada)} & \textbf{0,7417} & [0,733; 0,751] & \textbf{0,7946} & 0,8123 & 0,7798 & 0,8042 \\
Solo el texto (desarme interno)     & 0,6551 & [0,645; 0,665] & 0,7019 & 0,7191 & 0,7033 & 0,7092 \\
Cascada con el reparto anterior     & 0,5191 & [0,497; 0,542] & 0,5562 & 0,6601 & 0,5778 & 0,6446 \\
\bottomrule
\end{tabular}
\end{center}

\begin{itemize}\setlength{\itemsep}{0.1em}
  \item \textbf{25 de 30 familias} superan 0,50 y \textbf{21 de 30} superan 0,70.
  \item Mejora del reparto corregido: \textbf{+0,2225} [+0,199; +0,246], favorable en las 50 semillas.
  \item Aporte de las capas de reglas sobre el texto solo: \textbf{+0,0866}.
  \item La capa de reglas, aislada, \textbf{acierta 0,9928} y resuelve 80,3 de las 149 notas (53,9\,\%).
\end{itemize}

\subsection*{Sobre los intervalos}

Los de la tabla miden cuánto se mueve la cifra al cambiar la partición. Medimos además el intervalo
que corresponde para hablar del \textbf{corpus}, remuestreando \emph{plantillas} y no notas, porque
las notas de una misma plantilla son casi copias y tratarlas como observaciones independientes
estrecha artificialmente el intervalo.

\begin{center}
\small
\begin{tabular}{lcc}
\toprule
\textbf{Capa} & \textbf{Intervalo al 95\,\%} & \textbf{Escenario más conservador} \\
\midrule
\textbf{Sistema completo} & \textbf{[0,659; 0,819]} & [0,587; 0,794] \\
Solo el texto             & [0,565; 0,745] & [0,507; 0,718] \\
\bottomrule
\end{tabular}
\end{center}

La columna principal remuestrea las plantillas \emph{dentro de cada familia}, que es la pregunta que
corresponde al diseño: las treinta familias están fijadas por el conjunto de referencia, y lo que
podría haber salido distinto es qué plantillas conseguimos de cada una. La columna de la derecha
remuestrea sin distinguir familia, lo que además admitiría corpus sin alguna de ellas; la incluimos
porque es el escenario más desfavorable. \textbf{La conclusión es la misma en las dos: el intervalo
entero queda por encima del umbral}, y el sistema completo conserva más margen que el texto solo.

Este cálculo se hizo \textbf{dos veces, por separado}, con implementaciones escritas de forma
independiente y volviendo a calcular las predicciones desde cero en cada una. Las dos coinciden en
los cuatro decimales, en las dos columnas.

\subsection*{Por qué las reglas deciden antes que el texto}

No es una elección de comodidad: es la configuración medida como mejor. Si se hace que las reglas
\emph{cedan} ante el texto cuando el texto está muy seguro, el sistema empeora de forma sistemática
y monótona, y la degradación es mayor cuanto antes ceden.

\begin{center}
\small
\begin{tabular}{lcc}
\toprule
\textbf{Las reglas ceden\ldots} & \textbf{Cambio en macro-F1} & \textbf{Semillas a favor} \\
\midrule
siempre (equivale a texto solo) & $-0,0866$ & 0 de 50 \\
con confianza alta del texto    & $-0,0321$ & 0 de 50 \\
con confianza muy alta          & $-0,0120$ & 2 de 50 \\
casi nunca                      & $-0,0062$ & 2 de 50 \\
\bottomrule
\end{tabular}
\end{center}

Se exploró además una variante de cesión en el extremo, que da $+0,0015$. \textbf{No se adopta}: la
ganancia es del 0,2\,\% relativo, el punto favorable es estrecho, y el umbral se habría elegido
sobre el mismo conjunto con el que se mide. Queda documentada como explorada y descartada.

Hay además una verificación directa de por qué esas reglas funcionan. Se tomaron los seis pares de
familias que más texto comparten entre sí y se cruzaron sus marcadores --- correos de contacto,
servicios ocultos, monederos y enlaces. \textbf{Ninguno de los seis comparte un solo dato de
contacto.} El único valor que aparece en dos familias a la vez es el enlace de descarga del
navegador Tor, que no identifica a nadie, y que nuestro filtro de genéricos descarta
automáticamente por aparecer en más de una familia. Es decir: \textbf{aunque dos familias copien el
texto una de la otra, sus marcadores siguen siendo propios}, que es exactamente el supuesto sobre
el que se apoya la primera capa de la cascada.

\subsection*{Comportamiento en uso real}

El sistema puede \textbf{abstenerse} cuando la decisión del texto es dudosa. Las capas de reglas
siempre responden, porque su acierto ya es 0,99.

\begin{center}
\small
\begin{tabular}{lcc}
\toprule
\textbf{Umbral} & \textbf{Responde} & \textbf{Acierta donde responde} \\
\midrule
Sin abstención & 100\,\% & 0,8123 \\
0,50 & 77,2\,\% & \textbf{0,9324} \\
1,00 & 66,6\,\% & \textbf{0,9864} \\
\bottomrule
\end{tabular}
\end{center}

Y como apoyo a un analista, que es el uso previsto, la familia correcta es la \textbf{primera
propuesta en el 81,2\,\%} de las notas y está \textbf{entre las tres primeras en el 86,6\,\%}.

\section*{4. Dónde falla, y por qué}

\textbf{Las dos familias de plantilla única, BADRABBIT y CRYPTOLOCKER, dan F1 = 0.} Es estructural:
con una sola plantilla, evaluar sobre plantilla nunca vista implica no tener nada con que entrenar.
Ningún protocolo honesto lo evita, y lo declaramos así. Responde directamente a su observación de
que esas familias «no se procesan». Por eso informamos también sobre las \textbf{28 evaluables}.
Las otras tres por debajo de 0,50 son HELLOKITTY (0,293), RYUK (0,358) y JIGSAW (0,499).

\textbf{Errores de linaje.} Buena parte del error restante es confusión entre familias
emparentadas, no error arbitrario. La cascada baja los errores totales de 2.093 a 1.398, y los de
linaje de 486 a 186: en fracción, del 23,2\,\% al \textbf{13,3\,\%} de los errores. La confusión más
frecuente del texto es DHARMA con PHOBOS, dos familias de linaje común, con 8,1 casos por semilla.

Lo cuantificamos con un control: si se contabilizan como acierto las confusiones dentro de los pares
emparentados, el acierto del sistema sube de 0,8123 a 0,8585 sobre 27 clases. Como fusionar clases
sube el resultado por sí solo, comparamos contra fusionar la misma cantidad de pares \emph{elegidos
al azar} entre familias sin parentesco, promediado sobre doscientos sorteos: el azar solo llega a
0,8132. Es decir que \textbf{casi toda la ganancia es real}, y se traduce así: de los 18,8 puntos de
error del sistema, unos 4,6 son confundir dos familias que comparten el molde de la nota, cerca de
una cuarta parte del error total.

\section*{5. Limitaciones declaradas}

\textbf{1. El criterio de plantilla no detecta contención.} Agrupamos por similitud de caracteres,
que no ve cuando una nota está \emph{contenida} dentro de otra. Es la contracara del desglose de la
Sección~3: esas 54 notas en las que el sistema acierta 0,9891 cuentan como «plantilla nunca vista»
según nuestro criterio y, sin embargo, tienen una hermana muy parecida en el entrenamiento. Sin esa
ventaja el sistema acierta 0,7434, y el texto solo, 0,5839.

Toda cifra de este informe corresponde a «plantilla no vista \emph{según el criterio de
casi-duplicado declarado}». Con un criterio de contención más estricto, el corpus baja de 99 a 81
plantillas y de 28 a 21 familias evaluables, y las cifras bajan en consecuencia.

\textbf{2. Texto compartido entre familias.} Un barrido por $n$-gramas de ocho palabras detectó
notas que comparten texto con familias ajenas. Los seis pares con más texto en común \textbf{se
verificaron uno por uno abriendo los archivos, y ninguno es un error de etiquetado}. Lo que decide
si el texto compartido importa no es la cantidad sino la exclusividad: existe un molde de
ecosistema, frases que aparecen en seis familias distintas y cuyo uso compartido no dice nada. Los
pares verificados son además dos fenómenos distintos, y los nombramos distinto:

\begin{itemize}\setlength{\itemsep}{0.15em}
  \item \textbf{Un molde de nota compartido}, el caso de CLOP y RYUK: las dos notas empiezan con el
        mismo bloque continuo de cien palabras y comparten un prefijo literal de 423 caracteres.
        Provienen de repositorios públicos que ya las separan en dos familias, y cada una conserva
        sus propios contactos, su monedero y su firma.
  \item \textbf{Bloques de texto reutilizados}, el caso de los pares con SODINOKIBI y el de NOTPETYA
        con WANNACRY: comparten pasajes sueltos --- la advertencia de no usar software de
        recuperación, la enumeración de datos exfiltrados, la amenaza de publicación --- pero no el
        comienzo de la nota. Es práctica común del sector, no parentesco, y provienen de fuentes
        independientes entre sí.
\end{itemize}

\textbf{3. Redundancia del corpus público.} Las fuentes públicas de notas se copian entre sí. Es un
límite del material disponible, no del método, y está documentado nota por nota.

\textbf{4. Procedencia por auditar en un subconjunto.} Un grupo de notas quedó registrado con
procedencia agregada en vez de individual. Todas tienen fuente citable, pero la auditoría fina de
ese subconjunto sigue pendiente y lo declaramos.

\section*{6. Respuestas a sus pedidos}

\textbf{Métricas.} Además de macro-F1 informamos exactitud, exactitud balanceada, \textbf{coeficiente
de Matthews (0,8042)}, coeficiente de variación del F1 entre semillas (0,042), intervalos de
confianza al 95\,\%, F1 por familia y matriz de confusión.

\textbf{«¿Y las familias de una sola plantilla?»} Son dos, están identificadas, dan cero por
construcción, y por eso reportamos en paralelo sobre las 28 evaluables.

\textbf{Cuánto material hace falta por familia.} El aporte de sumar notas se agota pronto. Lo
medimos con deltas apareados bajo el mismo protocolo de la cifra de cabecera, sobre 50
repeticiones:

\begin{center}
\small
\begin{tabular}{lcc}
\toprule
\textbf{Notas por familia en entrenamiento} & \textbf{macro-F1} & \textbf{¿el salto aporta?} \\
\midrule
1 & 0,5757 & --- \\
2 & 0,6503 & sí \\
\textbf{3} & \textbf{0,6610} & \textbf{sí, y es el último que aporta} \\
4 & 0,6607 & no \\
6 y más & 0,660 y baja & no \\
\bottomrule
\end{tabular}
\end{center}

\textbf{A partir de la tercera nota por familia, sumar material no mejora el resultado.} La medición
equivalente contando \emph{textos distintos} en vez de notas da el mismo corte en tres, pero
proviene de un esquema de evaluación que entrena con un 42\,\% más de material que el de la cifra de
cabecera; lo señalamos porque los dos números no salen del mismo régimen. Lo común a los dos: el
corte está en tres y no se mueve, de modo que la conclusión no dependía del defecto de partición que
corregimos. Tiene una consecuencia práctica: \textbf{el límite de este frente no se levanta
recolectando más notas de las familias que ya tenemos, sino consiguiendo textos genuinamente
distintos de las que hoy tienen uno o dos}.

\textbf{Año de aparición frente a rendimiento.} No hay correlación significativa: $r = +0{,}091$ con
$p = 0{,}645$ sobre las 28 familias evaluables; descontando el número de plantillas, $r = +0{,}060$.
Es un análisis posterior y con 28 casos, así que se informa como exploratorio.

\section*{7. Aporte de este frente}

Además del sistema, la tesis aporta la \textbf{medición de cuánto del rendimiento reportado en esta
línea de trabajo es artefacto de evaluación}. Mostramos, con números, tres cosas: que casi la mitad
de las notas de un corpus público tiene una copia casi idéntica de sí misma; que evaluar sin separar
esas copias infla el resultado; y cuánto rendimiento queda al retirar cada fuente de optimismo. La
escalera completa de protocolos, desde el escenario de mundo cerrado hasta el más estricto, se
informa en el Capítulo~4.

\section*{8. Estado}

Las mediciones del frente de notas están \textbf{cerradas}. El Capítulo~4 ya tiene escritas la
cascada y el protocolo corregido, en seis subsecciones nuevas con sus tablas; solo se agregó, y las
cifras anteriores no se tocaron.

\vspace{0.6em}
\hrule
\vspace{0.5em}
{\footnotesize Todas las cifras de este informe están verificadas contra los archivos de resultados
y revisadas de forma cruzada por dos implementaciones independientes. Detalle, preregistros y
predicciones falladas en \texttt{ESTADO\_TESIS.md}.}

\end{document}
